\documentclass[11pt,a4paper]{article}
% Compile with XeLaTeX (fontspec and xeCJK).
\usepackage{fontspec}
\setmainfont{LibertinusSerif}[Extension=.otf,UprightFont=*-Regular,BoldFont=*-Bold,ItalicFont=*-Italic,BoldItalicFont=*-BoldItalic]
\setsansfont{LibertinusSans}[Extension=.otf,UprightFont=*-Regular,BoldFont=*-Bold,ItalicFont=*-Italic]
\usepackage{xeCJK}
\setCJKmainfont{Hiragino Mincho ProN}
\xeCJKDeclareCharClass{Default}{"00B7, "2014, "2018 -> "201F, "2026}
\usepackage[english]{babel}
\usepackage[a4paper,margin=2.5cm,top=2.3cm]{geometry}
\usepackage{graphicx}
\usepackage{xcolor}
\usepackage{titlesec}
\usepackage{fancyhdr}
\usepackage[hidelinks]{hyperref}
\definecolor{accent}{RGB}{31,73,125}
\definecolor{muted}{RGB}{93,104,119}
\hypersetup{colorlinks=true,urlcolor=accent,linkcolor=accent,pdfauthor={Romain Fernex},pdftitle={The origins of Korea and the kingdom of Goguryeo}}
\providecommand{\tightlist}{\setlength{\itemsep}{0pt}\setlength{\parskip}{0pt}}
\providecommand{\ul}[1]{\underline{#1}}
\setcounter{secnumdepth}{0}
\titleformat{\section}{\large\bfseries\color{accent}}{}{0em}{}
\setlength{\parindent}{1.2em}
\setlength{\parskip}{0.35em}
\pagestyle{fancy}
\fancyhf{}
\fancyhead[L]{\small\color{muted}Romain Fernex}
\fancyhead[R]{\small\color{muted}\itshape The origins of Korea and the kingdom of Goguryeo}
\fancyfoot[C]{\small\color{muted}\thepage}
\renewcommand{\headrulewidth}{0pt}
\newcommand{\articlefigure}[2]{%
  \begin{center}
  \includegraphics[width=0.55\linewidth,height=6.5cm,keepaspectratio]{#1}%
  \ifx\relax#2\relax\else\par\smallskip{\small\itshape\color{muted}#2\par}\fi
  \end{center}}

\begin{document}
\thispagestyle{empty}
\noindent{\small\sffamily\color{accent}\MakeUppercase{History of Korea · Part II}}\par\medskip
\noindent{\LARGE\bfseries The origins of Korea and the kingdom of Goguryeo\par}
\smallskip\noindent{\large\itshape Between nationalist narratives and historical realities\par}
\medskip
\noindent{\color{accent}\rule{\linewidth}{0.8pt}}\par\smallskip
\noindent{\small\color{muted}Romain Fernex · 2022 · Article in English}\par\bigskip

In the preceding part, we examined the period before the emergence of Goguryeo and the importance and nature of China's involvement in the development of nation-states on the peninsula. By looking at the populations living alongside one another there, we were able to highlight the facts and limits of the domination claimed by the People's Republic of China over this territory. Yet restricting ourselves to an analysis of this area's ethnic diversity does not fully address the geopolitical issues raised by the controversy over the \emph{Northeast Project}. Indeed, China's aim is to deny the political reality of a kingdom structured around a culture distinct from Chinese culture, laying the foundations of what would ultimately become the Korea we know. Likewise, for both South and North Korea, the aim is to root their existence and identity in ancient history by making this first kingdom of Goguryeo a cornerstone of their development.

It is therefore important to examine the kingdom itself and what made it distinctive, geopolitically, culturally and economically.

To do so, we will first go back again, although over much shorter timescales, to better understand the key moment of its emergence and how it developed before the Chinese conquest of 245, which made it ``disappear from history'' for nearly fifty years before its great return. We will then see how it fitted into and evolved within the ``Three Kingdoms'' order introduced at the end of the preceding article. Finally, we will examine the causes of its decline in favour of its rival, Silla, and try to answer once and for all the question raised in the preceding article.

Let us therefore return to 37 BC, the historical and mythological date of Goguryeo's foundation. Its founder is often known as Chumong, also called Dongmyeong, and is said to have originated in the kingdom of Buyeo, founded in the 2nd century BC in northern Manchuria. At its emergence, Goguryeo was more akin to a tribal confederation structured around five Bus{[}1{]}, which would soon become administrative divisions of the kingdom. These five Bus themselves originated in smaller tribal groupings called Nas, which still retained relative autonomy at the beginning of this period. As the Bus increased their power, one in particular, the Kyeru Bu, supplanted the others and would provide the kings of the dynasties ruling Goguryeo. This grouping was both the product and the cause of a policy of assimilation and conquest of the region's other chiefdoms and tribes, which would form its sphere of influence. Among the nomadic tribes that fell under Goguryeo's rule, several, including the Khitan and Malgal, retained substantial autonomy from the five Bus, although they remained subject to royal power. However, for several decades, the geographical and political unit within which Goguryeo's population lived and with which it interacted most was the eumnak. Eumnaks were communities of varying size, under village chiefs, through which the work and lives of the entire common population were organised.

\articlefigure{media/media/image1.png}{Map showing the main political entities of the Korean peninsula in 37 BC}

(the four commanderies of Han China are in purple)

Beyond these administrative considerations, an initial social hierarchy can already be observed. At the top of the social ladder were the Daega. This term could refer to the lord heading the Kyeru Bu, that is, the king, to the chiefs of the other Bus, or to local chiefs. Further down, the Sogas formed a property-owning petty bourgeoisie, generally receiving a salary from the government, not participating in agricultural activities, and able to aspire to a certain degree of social advancement. It should be noted that they could also be demoted to the lower class if they became too heavily indebted or disobeyed the court's orders. Most people belonged to the Haho category (tenants in Chinese), so called because its members owned no land. This class included peasants, artisans, merchants and others serving nobles but considered to be employed by the central government{[}2{]}, or directly under the control of great houses. The Haho class was far from homogeneous and included Homin, wealthy individuals who generally derived their fortunes from trade or crafts. Lower down the social ladder were the Yongjak, peasants who hired out their services to a wealthy landowner in exchange for a share of the harvest or a plot to cultivate. Finally, at the very bottom were the Nobi, servants directly employed by large landowners. These were often prisoners of war or peasants unable to repay their debts.

Economically, the royal authorities quickly introduced an initial system of taxation. Nevertheless, its effectiveness remained fairly limited, since it often had to operate through clan chiefs{[}3{]} and took the form of a simple tribute paid to the central authority.

Thus, almost from its emergence, Goguryeo society appeared as a complex system, different from that prevailing in China at the time and from the one China would establish after its first conquest of Goguryeo in 245. However, the effectiveness of the central authority claimed by the Daegas, and particularly the royal house, remained an open question. This problem was all the more acute because the kingdom's successive conquests had enabled agricultural advances and an expansion of cultivable land, fuelling Haho peasants' desire for individual ownership, although their labour was vital to the kingdom's wealth and functioning. Faced with peasants' growing demand for autonomy, a genuine centralising movement initiated by royal power enabled Goguryeo to emerge as a nation-state. This movement also arose from the need to develop an ambitious foreign policy, particularly in response to Han China, whose Liaodong commandery had harassed Goguryeo since its foundation and had already pushed it back to the Yalu valley. Far from a rapid transformation, it was the result of a slow, centuries-long process that began with the kingdom's creation and experienced several reversals throughout its history.

But how did this centralising movement unfold, and what evidence attests to its existence so early in Goguryeo's life?

An initial important element was the establishment of the political system of principalities. This centred on the king, who concentrated most power and ruled over the entire kingdom. Immediately beneath him were the chiefs of the five Bus forming the principalities, themselves divided into sub-principalities or Che-Bu. This first group of principalities covered most of the territory. Although four of these great principalities were subordinate to royal power, they retained an important consultative role and could sometimes exert real influence over the king's decisions. They coexisted with a second group comprising, on the one hand, other tribes accountable to the principalities (and especially to the king) but not directly attached to them (subordinate groups, see diagram below), and, on the other, chiefdoms that had fallen under Goguryeo's influence while retaining considerable autonomy in managing their territory and forming marquisates {[}4{]}.

\articlefigure{media/media/image2.png}{Diagram of Goguryeo's administrative organisation before AD 245}

From: Noh Taedon, \emph{Korea's Ancient Koguryo Kingdom: A Socio-political History}, Global Oriental, 2014, p.97

This transformation of Goguryeo's political system was accompanied by the gradual disappearance of eumnaks under pressure from Haho seeking property ownership and a government that saw community chiefs as an obstacle to its power, preventing it from exercising authority directly over peasants. Several instruments therefore developed to address this disconnection between royal power and the population, the source of its economic and military strength.

The first measure used by the monarchy to unify the population was the introduction of a cult of the king and his descendants based on the myth surrounding the birth of Chumong, Goguryeo's founder, who was said to be the son of the agricultural goddess Habaek Nyo. As in many monarchies, this was also a powerful means of establishing rulers' legitimacy, and it amply fulfilled that role until Goguryeo's fall.

The second was an initial reform of taxation, gradually shifting from a relatively rudimentary system of tribute demands to a genuine, systematic tax-collection system, making it a much better instrument of fiscal policy. However, this transformation extended over several centuries and occurred more substantially during the provincialisation of Goguryeo's political system, which we will discuss in the next article.

Finally, in international relations, Goguryeo distinguished itself as a remarkable military power, particularly through its cavalry and its mastery of ironworking, inherited in part from Buyeo, which had made this its speciality. Foreign trade also developed with the commanderies of the Wei dynasty, but these relations remained marked by mistrust, were often limited to tributes paid by China to certain tribes making up Goguryeo, frequently to encourage internal divisions, and could hardly obscure the climate of tension between the two kingdoms. These tensions culminated

in 245, when the Wei besieged and defeated Goguryeo's capital in retaliation for its policy of alliances with Chinese dynasties rival to the Wei.

\articlefigure{media/media/image3.png}{Map of East Asia in AD 262 showing the most influential Chinese dynasties}

The 3rd century thus ended in a dark age for Goguryeo, but one that had seen the foundations of a nation-state emerge, particularly politically and, admittedly less markedly, economically. We also observe the beginnings of cultural unity through the cult devoted to the royal family, although it remained at a primitive stage. In the next article, we will therefore analyse how this regime was reborn barely fifty years after its destruction, how it continued what had emerged since its foundation, and how it departed from it.

Romain Fernex, dual-degree student at ESCP and Korea University

\textbf{Bibliography:}

\begin{itemize}
\item
  Noh Taedon, \emph{Korea's Ancient Koguryo Kingdom : A Socio-political History}, global oriental, 2014
\item
  Kim Chang-seok, \emph{Goguryeo Society and Its Economy}, The Review of Korean Studies, Vol.7 n°4, 2004
\item
  Michael J.Seth,~\emph{A history of Korea -- From Antiquity to the present}, 2010, Rowman \& Littlefield Publishers
\item
  Jung Woon Yong,\emph{~Trends in Koguryo's relationship with Paekche and Silla during the 4th -- 7th centuries}, International Journal of Korean history, Vol. 8, 2006
\item
  Jeon Ho-tae\emph{, Goguryeo Culture and Mural Paintings}, The Review of Korean Studies, Vol.7 n°4, 2004
\end{itemize}

\textbf{Notes:}

\textbf{{[}1{]}} Although the term Bu as an administrative division of Goguryeo is uncontroversial, the same cannot be said of the interpretation of its meaning or meanings in other contexts. When it designates the entity preceding these administrative divisions, it is particularly difficult to define. Some historians initially suggested that a Bu in its earliest form was a tribal grouping, but the term tribe can sometimes pose a problem. Indeed, following L. H. Morgan's definition in \emph{Ancient Society} in 1877, a tribe can be defined as an intermediate state between a classless, disorganised society and a nation-state. Yet a relatively developed hierarchy can be observed very early within the Bus, potentially calling this initial definition of a Bu as a tribal grouping into question.

\textbf{{[}2{]}} When I say that some Haho were employed by the central government, this clearly does not mean a direct relationship of exploitation, and it often operated through great local families. However, the major difference from the second category of Haho introduced immediately afterwards is that they served the interests of the state rather than private interests.

\textbf{{[}3{]}} This is all the more true of the royal authorities' relations with the nomadic tribes subject to them but retaining considerable autonomy, such as the Khitans or the Malgals. In these particular cases, the tribute collected was much smaller than for the eumnaks and always passed through the clan chiefs, who thus acted as intermediaries. However, this autonomy was offset by the fact that the state could mobilise them at will for military or economic purposes whenever the situation required it.

\textbf{{[}4{]}} The term is used by Noh Taedon in \emph{Korea's Ancient Koguryo Kingdom: A Socio-political History}, but he explains that, although it has no basis in the Korean literature of the time, it is the most appropriate term here, given the strong similarities in functioning between the ``marquisates'' mentioned here and their European counterparts, which appeared much later.

\end{document}
